% ===========================================================================
%  Chapitre 2 — Systèmes d'équations linéaires
%  PE 2026, composante ALGÈBRE — RAPE, première période
% ===========================================================================
\bmtchapitre{Systèmes d'équations linéaires}
  {Résoudre dans $\R^{2}$ un système de deux équations à deux inconnues et
   dans $\R^{3}$ un système de trois équations à trois inconnues, par la
   méthode de Cramer comme par substitution, et mettre en équation un
   problème de la vie courante.}
  {Algèbre \textbullet\ environ 6 heures}

Un système linéaire, c'est plusieurs contraintes qui doivent tenir ensemble.
« Trois kilos de riz et deux litres d'huile coûtent tant ; deux kilos de riz
et cinq litres d'huile coûtent tant » : deux phrases, deux inconnues, une
seule réponse possible. Tout le chapitre consiste à passer de ces phrases à
deux lignes de calcul.

Deux méthodes coexistent, et il faut savoir les deux. La substitution est
naturelle et fonctionne toujours, mais elle s'alourdit vite à trois
inconnues. La méthode de Cramer est mécanique : on calcule des déterminants,
on divise, c'est fini. Le programme officiel la nomme explicitement pour le
système à trois inconnues.

% ---------------------------------------------------------------------------
\begin{revision}
Deux notions de seconde, rien de plus.

\begin{enumerate}[label=\textbf{\arabic*.}, leftmargin=*, itemsep=0.35em]
  \item Résoudre $5x-2 = 3x+8$.
  \item Exprimer $y$ en fonction de $x$ dans l'égalité $2x+3y = 12$.
  \item Deux droites du plan : combien de points communs peuvent-elles
        avoir ?
  \item Donner l'équation réduite de la droite passant par $A(0\,;\,3)$ et
        de coefficient directeur $-2$.
  \item Calculer $3\times7-4\times5$.
  \item Le point $B(1\,;\,-1)$ appartient-il à la droite d'équation
        $y = 2x-3$ ?
\end{enumerate}
\end{revision}

\begin{automatismes}
Sans calculatrice.

\begin{enumerate}[label=\textbf{\alph*)}, leftmargin=*, itemsep=0.25em]
  \item $2\times5-3\times1$
  \item $(-2)\times4-3\times(-1)$
  \item Si $x+y = 10$ et $x-y = 2$, que vaut $x$ ?
  \item Si $y = 2x$ et $x+y = 9$, que vaut $x$ ?
  \item $\dfrac{-12}{-4}$
  \item $\dfrac{0}{7}$
  \item Résoudre $4x = 0$
  \item Résoudre $0\times x = 5$
\end{enumerate}
\end{automatismes}

\begin{corrige}[automatismes]
a) $7$ \quad b) $-5$ \quad c) $6$ \quad d) $3$ \quad e) $3$ \quad f) $0$
\quad g) $x = 0$ \quad h) aucune solution.
\end{corrige}

% ===========================================================================
\section{Le système de deux équations à deux inconnues}

\begin{definition}[système linéaire $2\times2$]
On appelle \textbf{système linéaire de deux équations à deux inconnues} tout
système de la forme
\[
  (S) \quad
  \begin{cases}
    ax+by = c\\
    a'x+b'y = c'
  \end{cases}
\]
où $a$, $b$, $c$, $a'$, $b'$, $c'$ sont des réels donnés. Résoudre $(S)$,
c'est déterminer tous les couples $(x\,;\,y)$ de $\R^{2}$ qui vérifient
\emph{les deux} égalités à la fois.
\end{definition}

\begin{visualisation}[deux droites, trois situations]
\begin{center}
\begin{tikzpicture}
\begin{axis}[
  name = A, width = 4.9cm, height = 4.3cm, axis lines = middle,
  xmin = -2.4, xmax = 2.4, ymin = -2.2, ymax = 2.4,
  xtick = \empty, ytick = \empty, axis line style = {bmtGris},
  title = {\footnotesize un point commun}, title style = {color=bmtGris},
]
  \addplot[bmtIndigo, line width=1.1pt, domain=-2:2] {0.8*x+0.4};
  \addplot[bmtLaterite, line width=1.1pt, domain=-2:2] {-1.1*x+0.1};
  \filldraw[bmtVetiver] (axis cs:-0.158,0.274) circle (1.9pt);
\end{axis}
\begin{axis}[
  name = B, at = {($(A.east)+(1.15cm,0)$)}, anchor = west,
  width = 4.9cm, height = 4.3cm, axis lines = middle,
  xmin = -2.4, xmax = 2.4, ymin = -2.2, ymax = 2.4,
  xtick = \empty, ytick = \empty, axis line style = {bmtGris},
  title = {\footnotesize aucun point commun}, title style = {color=bmtGris},
]
  \addplot[bmtIndigo, line width=1.1pt, domain=-2:2] {0.8*x+0.6};
  \addplot[bmtLaterite, line width=1.1pt, domain=-2:2] {0.8*x-0.7};
\end{axis}
\begin{axis}[
  name = C, at = {($(B.east)+(1.15cm,0)$)}, anchor = west,
  width = 4.9cm, height = 4.3cm, axis lines = middle,
  xmin = -2.4, xmax = 2.4, ymin = -2.2, ymax = 2.4,
  xtick = \empty, ytick = \empty, axis line style = {bmtGris},
  title = {\footnotesize tous les points}, title style = {color=bmtGris},
]
  \addplot[bmtIndigo, line width=2.6pt, domain=-2:2, opacity=0.35] {0.8*x+0.2};
  \addplot[bmtLaterite, line width=1.1pt, domain=-2:2] {0.8*x+0.2};
\end{axis}
\end{tikzpicture}
\end{center}

Chaque équation du système est celle d'une droite. Résoudre le système, c'est
donc chercher les points communs à ces deux droites — et il n'y a que trois
possibilités : elles se coupent en un point, elles sont strictement
parallèles, ou elles sont confondues. C'est exactement ce que le déterminant
va dire, sans qu'on ait à tracer quoi que ce soit.
\end{visualisation}

\begin{definition}[déterminant d'ordre 2]
Le \textbf{déterminant} du système $(S)$ est le nombre
\[
  D = \begin{vmatrix} a & b\\ a' & b'\end{vmatrix} = ab'-a'b .
\]
\end{definition}

\begin{theoreme}[méthode de Cramer, ordre 2]
Posons
\[
  D_{x} = \begin{vmatrix} c & b\\ c' & b'\end{vmatrix} = cb'-c'b,
  \qquad
  D_{y} = \begin{vmatrix} a & c\\ a' & c'\end{vmatrix} = ac'-a'c .
\]
Si $D \neq 0$, le système admet un unique couple solution :
\[
  x = \frac{D_{x}}{D}, \qquad y = \frac{D_{y}}{D}.
\]
Si $D = 0$, le système n'a soit aucune solution, soit une infinité.
\end{theoreme}

\begin{demonstration}
Multiplions la première équation par $b'$, la seconde par $b$, puis
soustrayons : les termes en $y$ disparaissent et il reste
$(ab'-a'b)\,x = cb'-c'b$, c'est-à-dire $Dx = D_{x}$. Le même procédé, en
éliminant $x$, donne $Dy = D_{y}$. Si $D \neq 0$, on divise, et l'on vérifie
réciproquement que le couple obtenu convient.
\end{demonstration}

\begin{remarque}
Retenir la place des colonnes plutôt que les formules : dans $D_{x}$, la
colonne des $x$ est remplacée par celle des seconds membres ; dans $D_{y}$,
c'est la colonne des $y$. Ce principe vaut à l'identique pour les systèmes à
trois inconnues.
\end{remarque}

\begin{exemple}[le même système, deux méthodes]
\[
  \begin{cases} 3x+2y = 16\\ x-y = 2 \end{cases}
\]

\emph{Par substitution.} La seconde équation donne $y = x-2$. En reportant
dans la première : $3x+2(x-2) = 16$, soit $5x = 20$ et $x = 4$, d'où
$y = 2$.

\emph{Par Cramer.} $D = 3\times(-1)-1\times2 = -5$,
$D_{x} = 16\times(-1)-2\times2 = -20$,
$D_{y} = 3\times2-1\times16 = -10$. D'où
$x = \frac{-20}{-5} = 4$ et $y = \frac{-10}{-5} = 2$.

Les deux méthodes donnent le même couple $(4\,;\,2)$, et l'on vérifie :
$3\times4+2\times2 = 16$. \checkmark
\end{exemple}

\begin{methode}{choisir sa méthode}
\begin{itemize}[leftmargin=1.3em, itemsep=0.2em]
  \item Une inconnue est déjà isolée, ou son coefficient vaut $1$ ou $-1$ :
        \textbf{substitution}.
  \item Les coefficients sont quelconques, ou l'on veut un calcul sans
        réflexion : \textbf{Cramer}.
  \item On demande de discuter selon un paramètre : \textbf{Cramer}, car le
        déterminant répond seul à la question de l'existence.
\end{itemize}
\end{methode}

\begin{entrainement}[systèmes $2\times2$]
\exo[\niveau{1}] $\begin{cases} x+y = 7\\ x-y = 1\end{cases}$

\exo[\niveau{1}] $\begin{cases} 2x+3y = 8\\ x-y = -1\end{cases}$

\exo[\niveau{2}] $\begin{cases} 4x-3y = 6\\ 5x+2y = 19\end{cases}$
\quad (par Cramer)

\exo[\niveau{2}] $\begin{cases} 2x-y = 3\\ -6x+3y = 1\end{cases}$
\quad (calculer $D$ d'abord)

\exo[\niveau{2}] $\begin{cases} x+2y = 5\\ 3x+6y = 15\end{cases}$

\exo[\niveau{3}] Déterminer $m$ pour que le système
$\begin{cases} mx+y = 3\\ 4x+my = 6\end{cases}$
n'admette pas une unique solution.
\end{entrainement}

\begin{corrige}[systèmes $2\times2$]
\textbf{1.} $(4\,;\,3)$.
\textbf{2.} $y = x+1$, d'où $2x+3x+3 = 8$ et $x = 1$ : $(1\,;\,2)$.
\textbf{3.} $D = 8+15 = 23$, $D_{x} = 12+57 = 69$, $D_{y} = 76-30 = 46$ :
$(3\,;\,2)$.
\textbf{4.} $D = 6-6 = 0$. La seconde équation vaut $-3$ fois la première
côté gauche, mais $-3\times3 = -9 \neq 1$ : aucune solution.
\textbf{5.} $D = 0$, et la seconde équation est le triple de la première :
une infinité de solutions, tous les couples $\left(x\,;\,\frac{5-x}{2}\right)$.
\textbf{6.} $D = m^{2}-4$ ; le système n'a pas de solution unique lorsque
$D = 0$, soit $m = -2$ ou $m = 2$.
\end{corrige}

% ===========================================================================
\section{Le système de trois équations à trois inconnues}

\begin{definition}[système linéaire $3\times3$]
\[
  (S) \quad
  \begin{cases}
    ax+by+cz = d\\
    a'x+b'y+c'z = d'\\
    a''x+b''y+c''z = d''
  \end{cases}
\]
Une solution est un \textbf{triplet} $(x\,;\,y\,;\,z)$ vérifiant les trois
égalités.
\end{definition}

\begin{definition}[déterminant d'ordre 3, règle de Sarrus]
On recopie les deux premières colonnes à droite du tableau, puis on ajoute
les trois produits « descendants » et l'on retranche les trois produits
« montants » :
\[
  D = \begin{vmatrix}
    a & b & c\\ a' & b' & c'\\ a'' & b'' & c''
  \end{vmatrix}
  = ab'c''+bc'a''+ca'b''-ca'' b'-ab''c'-bc''a' .
\]
\end{definition}

\begin{visualisation}[la règle de Sarrus, en image]
\begin{center}
\begin{tikzpicture}[scale=1]
  \foreach \i/\j/\v in {0/0/a, 1/0/b, 2/0/c, 3/0/a, 4/0/b,
                        0/-1/{a'}, 1/-1/{b'}, 2/-1/{c'}, 3/-1/{a'}, 4/-1/{b'},
                        0/-2/{a''}, 1/-2/{b''}, 2/-2/{c''}, 3/-2/{a''}, 4/-2/{b''}}
    \node[font=\small] at (\i*1.25, \j*0.85) {$\v$};
  \draw[bmtGrisClair, line width=0.6pt] (2.65,0.45) -- (2.65,-2.15);
  \foreach \s in {0,1,2}
    \draw[bmtIndigo, line width=0.9pt, opacity=0.75]
      (\s*1.25-0.28, 0.28) -- (\s*1.25+2.22, -1.98);
  \foreach \s in {0,1,2}
    \draw[bmtLaterite, line width=0.9pt, opacity=0.75, dash pattern=on 3pt off 2pt]
      (\s*1.25-0.28, -1.98) -- (\s*1.25+2.22, 0.28);
  \node[bmtIndigo, font=\footnotesize, anchor=west] at (5.55,-0.2)
    {$+$ les trois descendantes};
  \node[bmtLaterite, font=\footnotesize, anchor=west] at (5.55,-1.35)
    {$-$ les trois montantes};
\end{tikzpicture}
\end{center}

La règle de Sarrus ne vaut que pour les déterminants d'ordre trois. Elle est
sûre à une condition : recopier vraiment les deux premières colonnes, au
lieu de tenter de les imaginer.
\end{visualisation}

\begin{theoreme}[méthode de Cramer, ordre 3]
Notons $D_{x}$, $D_{y}$, $D_{z}$ les déterminants obtenus en remplaçant,
dans $D$, la colonne des $x$, des $y$ ou des $z$ par celle des seconds
membres. Si $D \neq 0$, le système admet un unique triplet solution :
\[
  x = \frac{D_{x}}{D}, \qquad y = \frac{D_{y}}{D}, \qquad
  z = \frac{D_{z}}{D}.
\]
\end{theoreme}

\begin{exemple}[un système $3\times3$ résolu par Cramer]
\[
  \begin{cases}
    x+y+z = 6\\
    2x-y+z = 3\\
    x+2y-z = 2
  \end{cases}
\]

\emph{Le déterminant.} Par la règle de Sarrus, les trois produits descendants
donnent
\[
  (1)(-1)(-1)+(1)(1)(1)+(1)(2)(2) = 1+1+4 = 6,
\]
et les trois produits montants
\[
  (1)(1)(-1)+(1)(2)(1)+(1)(-1)(2) = -1+2-2 = -1 .
\]
D'où $D = 6-(-1) = 7$.
Comme $D \neq 0$, la solution est unique.

\emph{Les trois autres.} En remplaçant la première colonne par
$(6\,;\,3\,;\,2)$ on obtient $D_{x} = 7$ ; la deuxième colonne donne
$D_{y} = 14$ ; la troisième, $D_{z} = 21$.

\emph{Conclusion.} $x = 1$, $y = 2$, $z = 3$. On vérifie sur la première
équation : $1+2+3 = 6$. \checkmark
\end{exemple}

\begin{methode}{résoudre un système $3\times3$ par substitution (Gauss)}
Quand les coefficients s'y prêtent, on peut préférer éliminer les inconnues
l'une après l'autre.
\begin{enumerate}[label=\textbf{\arabic*.}, leftmargin=*, itemsep=0.15em]
  \item Choisir une équation où une inconnue a le coefficient $1$ ou $-1$ et
        l'isoler.
  \item Reporter dans les deux autres : on obtient un système $2\times2$.
  \item Le résoudre, puis remonter pour trouver la troisième inconnue.
  \item Vérifier le triplet dans l'équation qui n'a pas servi au report.
\end{enumerate}
\end{methode}

\begin{erreurclassique}
Dans un déterminant, l'ordre des lignes doit être le même que celui des
équations, et l'ordre des colonnes le même que celui des inconnues. Un
système où une équation s'écrirait $2z+x = 5$ doit d'abord être réordonné en
$x+0\times y+2z = 5$ — sans oublier le zéro. La moitié des erreurs de calcul
sur les systèmes $3\times3$ viennent d'un coefficient nul qu'on a « sauté ».
\end{erreurclassique}

\begin{entrainement}[systèmes $3\times3$]
\exo[\niveau{1}] $\begin{cases} x+y+z = 3\\ x-y+z = 1\\ x+y-z = 1\end{cases}$

\exo[\niveau{2}] $\begin{cases} 2x+y-z = 3\\ x-y+2z = 1\\ 3x+2y+z = 8\end{cases}$

\exo[\niveau{2}] $\begin{cases} x+2y+3z = 14\\ 2x-y+z = 3\\ 3x+y-z = 2\end{cases}$
\quad (par Cramer)

\exo[\niveau{3}] Calculer le déterminant
$\begin{vmatrix} 1 & 2 & 3\\ 2 & 4 & 6\\ 1 & 0 & 5\end{vmatrix}$
et dire ce que l'on peut en conclure sur le système correspondant.
\end{entrainement}

\begin{corrige}[systèmes $3\times3$]
\textbf{1.} En additionnant les deux premières : $2x+2z = 4$. En additionnant
la première et la troisième : $2x+2y = 4$. Avec $x+y+z = 3$, on trouve
$(1\,;\,1\,;\,1)$.
\textbf{2.} La première donne $z = 2x+y-3$ ; en reportant dans les deux
autres : $x-y+2(2x+y-3) = 1$, soit $5x+y = 7$, et
$3x+2y+2x+y-3 = 8$, soit $5x+3y = 11$. En soustrayant, $2y = 4$ donc
$y = 2$, puis $x = 1$ et $z = 1$. Le triplet est $(1\,;\,2\,;\,1)$.
\textbf{3.} Par la règle de Sarrus, $D = 25$. On obtient de même
$D_{x} = 25$, $D_{y} = 50$ et $D_{z} = 75$, d'où le triplet
$(1\,;\,2\,;\,3)$.
\textbf{4.} La deuxième ligne est le double de la première : le déterminant
est nul. Le système n'admet donc pas de solution unique — soit aucune, soit
une infinité, selon les seconds membres.
\end{corrige}

\begin{qcm}
\exo Le déterminant $\begin{vmatrix} 2 & -1\\ 3 & 4\end{vmatrix}$ vaut :
\textbf{a)} $5$ \quad \textbf{b)} $11$ \quad \textbf{c)} $-11$ \quad
\textbf{d)} $8$ \reponseqcm{b}

\exo Si le déterminant d'un système $2\times2$ est nul, le système :
\textbf{a)} a une solution unique \quad \textbf{b)} n'a jamais de solution
\quad \textbf{c)} a zéro ou une infinité de solutions \quad
\textbf{d)} a exactement deux solutions \reponseqcm{c}

\exo Dans la méthode de Cramer, $D_{y}$ s'obtient en remplaçant :
\textbf{a)} la ligne des $y$ \quad \textbf{b)} la colonne des $y$ par celle
des seconds membres \quad \textbf{c)} la colonne des seconds membres par
celle des $y$ \quad \textbf{d)} la première colonne \reponseqcm{b}

\exo La règle de Sarrus s'applique aux déterminants :
\textbf{a)} d'ordre 2 seulement \quad \textbf{b)} d'ordre 3 seulement \quad
\textbf{c)} de tout ordre \quad \textbf{d)} d'ordre pair \reponseqcm{b}
\end{qcm}

% ===========================================================================
\section{Mettre un problème en équations}

\begin{methode}{de la phrase au système}
\begin{enumerate}[label=\textbf{\arabic*.}, leftmargin=*, itemsep=0.15em]
  \item \emph{Nommer.} Écrire noir sur blanc ce que désignent $x$, $y$, $z$,
        \textbf{avec leur unité}. Cette ligne rapporte des points et évite la
        moitié des erreurs.
  \item \emph{Traduire.} Une phrase, une équation. « Au total », « en tout »
        donnent une somme ; « deux fois plus que » donne un produit.
  \item \emph{Résoudre} le système obtenu.
  \item \emph{Revenir à la question.} La réponse est une phrase, avec les
        unités, et l'on vérifie qu'elle a un sens : un effectif est entier,
        un prix est positif.
\end{enumerate}
\end{methode}

\begin{exemple}[trois achats au marché]
Au marché, trois kilos de riz et deux litres d'huile coûtent
$29\,000$ ariary ; deux kilos de riz et cinq litres d'huile coûtent
$50\,500$ ariary. Quel est le prix du kilo de riz ?

Notons $x$ le prix d'un kilo de riz et $y$ celui d'un litre d'huile, en
ariary. Le texte donne
\[
  \begin{cases} 3x+2y = 29\,000\\ 2x+5y = 50\,500 \end{cases}
\]
Ici $D = 3\times5-2\times2 = 11$. Pour $D_{x}$, on remplace la colonne des
$x$ par celle des seconds membres :
$D_{x} = 29\,000\times5-50\,500\times2 = 145\,000-101\,000 = 44\,000$, donc
$x = \frac{44\,000}{11} = 4\,000$. De même
$D_{y} = 3\times50\,500-2\times29\,000 = 151\,500-58\,000 = 93\,500$, donc
$y = \frac{93\,500}{11} = 8\,500$.

Le kilo de riz coûte $4\,000$ ariary et le litre d'huile $8\,500$ ariary. On
vérifie sur la seconde équation :
$2\times4\,000+5\times8\,500 = 8\,000+42\,500 = 50\,500$. \checkmark
\end{exemple}

\begin{entrainement}[mise en équation]
\exo[\niveau{1}] La somme de deux nombres est $48$ et leur différence $12$.
Quels sont ces nombres ?

\exo[\niveau{2}] Dans une classe de $42$ élèves, il y a huit filles de plus
que de garçons. Combien de filles ?

\exo[\niveau{2}] Un commerçant vend des sacs de $25$ kg et de $50$ kg. Il en
a vendu $23$ en tout, pour une masse totale de $850$ kg. Combien de sacs de
chaque sorte ?

\exo[\niveau{3}] Trois nombres ont pour somme $60$. Le deuxième est le double
du premier, et le troisième dépasse le premier de $10$. Déterminer ces trois
nombres.
\end{entrainement}

\begin{corrige}[mise en équation]
\textbf{1.} $x+y = 48$ et $x-y = 12$ : $30$ et $18$.
\textbf{2.} $f+g = 42$ et $f = g+8$ : $25$ filles et $17$ garçons.
\textbf{3.} $x+y = 23$ et $25x+50y = 850$, soit $x+2y = 34$ : $y = 11$ sacs
de $50$ kg et $x = 12$ sacs de $25$ kg.
\textbf{4.} $x+y+z = 60$, $y = 2x$, $z = x+10$ : $4x+10 = 60$, donc $x = 12{,}5$,
$y = 25$ et $z = 22{,}5$.
\end{corrige}

\begin{vraifaux}
\exo Un système de deux équations à deux inconnues a toujours une solution.

\exo Si $D = 0$, le système n'a aucune solution.

\exo Un système dont les deux équations représentent la même droite a une
infinité de solutions.

\exo La méthode de Cramer donne la solution même lorsque $D = 0$.

\exo Un système $3\times3$ dont deux équations sont identiques ne peut pas
avoir de solution unique.
\end{vraifaux}

\begin{corrige}[vrai ou faux]
\textbf{1.} Faux — deux droites strictement parallèles n'ont aucun point
commun.
\textbf{2.} Faux — il peut aussi en avoir une infinité ; $D = 0$ dit
seulement que la solution n'est pas unique.
\textbf{3.} Vrai — tous les points de cette droite conviennent.
\textbf{4.} Faux — on diviserait par zéro. Quand $D = 0$, il faut revenir aux
équations elles-mêmes.
\textbf{5.} Vrai — le déterminant est alors nul, deux lignes étant égales.
\end{corrige}

% ===========================================================================
\begin{situation}[la coopérative de vanille]
Une coopérative conditionne sa vanille en trois qualités :
extra, standard et courante. Une caisse mélangée contient $x$ kilos d'extra,
$y$ kilos de standard et $z$ kilos de courante. On sait que :

\begin{itemize}[leftmargin=1.3em, itemsep=0.15em]
  \item la caisse pèse $60$ kilos en tout ;
  \item elle contient deux fois plus de standard que d'extra ;
  \item l'extra se vend $90\,000$ ariary le kilo, le standard $60\,000$, le
        courant $30\,000$, et la caisse entière vaut $2\,700\,000$ ariary.
\end{itemize}

\begin{enumerate}[label=\textbf{\arabic*.}, leftmargin=*, itemsep=0.3em]
  \item Traduire ces trois renseignements par un système de trois équations à
        trois inconnues.
  \item Résoudre ce système.
  \item La coopérative peut-elle composer, avec les mêmes prix et la même
        masse, une caisse valant $3\,000\,000$ d'ariary contenant toujours
        deux fois plus de standard que d'extra ? Justifier.
\end{enumerate}
\end{situation}

\begin{corrige}[la coopérative de vanille]
\textbf{1.} $x+y+z = 60$ ; $y = 2x$ ; $90\,000x+60\,000y+30\,000z
= 2\,700\,000$, que l'on simplifie en $3x+2y+z = 90$.

\textbf{2.} En remplaçant $y$ par $2x$ : $3x+z = 60$ et $7x+z = 90$. En
soustrayant, $4x = 30$, donc $x = 7{,}5$, $y = 15$ et $z = 37{,}5$. La caisse
contient $7{,}5$ kg d'extra, $15$ kg de standard et $37{,}5$ kg de courante.

\textbf{3.} Le système devient $3x+z = 60$ et $7x+z = 100$, d'où $x = 10$,
$y = 20$ et $z = 30$. C'est possible : il suffit d'augmenter la part d'extra.
\end{corrige}

% ===========================================================================
\begin{annales}
Aucun sujet de baccalauréat de la série, parmi ceux dépouillés pour ce livre
— de la session 1999 à la session 2026 —, ne consacre d'exercice entier aux
systèmes linéaires. Ils y apparaissent en revanche systématiquement
\emph{à l'intérieur} de l'exercice de statistique : déterminer l'équation
d'une droite passant par deux points, c'est résoudre un système $2\times2$.
Les énoncés ci-dessous en sont les extraits, ramenés à leur seul contenu
algébrique — les coordonnées données sont celles calculées dans le sujet.

\exoann{Baccalauréat, Madagascar, série A, session 2026 — exercice 2,
question 3.b) (extrait)}
Déterminer les réels $a$ et $b$ tels que la droite d'équation $y = ax+b$
passe par les points $G_{1}(2\,;\,325)$ et $G_{2}(5\,;\,380)$.

\exoann{Baccalauréat, Madagascar, série A, session 2017 — exercice 2,
question 2.c) (extrait)}
Déterminer l'équation cartésienne de la droite passant par
$G_{1}(2\,;\,34)$ et $G_{2}(5\,;\,44)$.

\exoann{Baccalauréat, Madagascar, série A, session 2014 — exercice 2,
question 3 (extrait)}
La droite d'ajustement passe par $G_{1}(2\,;\,66)$ et $G_{2}(5\,;\,75)$.
Déterminer son équation, puis calculer l'ordonnée du point de cette droite
d'abscisse $11$.

\exoann{Baccalauréat, Madagascar, série A, session 2010 — exercice 1,
question 2.c) (extrait)}
Déterminer une équation de la droite $(G_{1}G_{2})$, où
$G_{1}(5{,}8\,;\,8{,}2)$ et $G_{2}(12{,}2\,;\,15)$.
\end{annales}

\begin{corrige}[annales]
Dans les quatre cas, on écrit que chacun des deux points vérifie
$y = ax+b$, ce qui donne un système $2\times2$ d'inconnues $a$ et $b$.

\textbf{1.} $\begin{cases} 2a+b = 325\\ 5a+b = 380\end{cases}$ : en
soustrayant, $3a = 55$, donc $a = \frac{55}{3}$ et
$b = 325-\frac{110}{3} = \frac{865}{3}$. La droite a pour équation
$y = \frac{55}{3}x+\frac{865}{3}$, soit environ $y = 18{,}33x+288{,}33$.

\textbf{2.} $3a = 10$, donc $a = \frac{10}{3}$ et $b = 34-\frac{20}{3}
= \frac{82}{3}$ : $y = \frac{10}{3}x+\frac{82}{3}$.

\textbf{3.} $3a = 9$, donc $a = 3$ et $b = 60$ : $y = 3x+60$. Pour $x = 11$,
$y = 93$.

\textbf{4.} $6{,}4\,a = 6{,}8$, donc $a = \frac{17}{16} = 1{,}0625$ et
$b = 8{,}2-1{,}0625\times5{,}8 = 2{,}0375$. La droite a pour équation
$y = 1{,}0625\,x+2{,}0375$.
\end{corrige}

\begin{remarque}[pourquoi ces extraits comptent]
Ce ne sont pas des exercices artificiels : la question « déterminer
l'équation de la droite $(G_{1}G_{2})$ » vaut un point entier à chaque
session, soit un vingtième de la note finale, et elle se résout exactement
comme les systèmes de ce chapitre. Vous la retrouverez complète au
chapitre 14.
\end{remarque}

% ===========================================================================
\begin{jecherche}[un système à paramètre]
On considère, pour tout réel $m$, le système
\[
  (S_{m}) \quad
  \begin{cases}
    x+my = 2\\
    mx+4y = m+2
  \end{cases}
\]

\begin{enumerate}[label=\textbf{\arabic*.}, leftmargin=*, itemsep=0.3em]
  \item Calculer le déterminant $D$ de $(S_{m})$ en fonction de $m$.
  \item Pour quelles valeurs de $m$ le système admet-il une solution unique ?
        La déterminer.
  \item Étudier séparément les cas $m = 2$ et $m = -2$.
\end{enumerate}
\end{jecherche}

\begin{corrige}[un système à paramètre]
\textbf{1.} $D = 4-m^{2}$.

\textbf{2.} $D \neq 0 \iff m \neq -2$ et $m \neq 2$. On calcule
$D_{x} = 2\times4-(m+2)m = 8-m^{2}-2m$ et
$D_{y} = (m+2)-2m = 2-m$. Donc
$y = \dfrac{2-m}{4-m^{2}} = \dfrac{2-m}{(2-m)(2+m)} = \dfrac{1}{m+2}$, et
$x = 2-my = 2-\dfrac{m}{m+2} = \dfrac{m+4}{m+2}$.

\textbf{3.} Pour $m = 2$, le système devient $x+2y = 2$ et $2x+4y = 4$ : la
seconde équation est le double de la première, il y a une infinité de
solutions, les couples $\left(2-2y\,;\,y\right)$.
Pour $m = -2$ : $x-2y = 2$ et $-2x+4y = 0$, soit $x = 2y$ ; en reportant,
$2y-2y = 2$, c'est-à-dire $0 = 2$ : aucune solution.
\end{corrige}

% ===========================================================================
\begin{jemeteste}
Vingt minutes.

\begin{enumerate}[label=\textbf{\arabic*.}, leftmargin=*, itemsep=0.3em]
  \item Résoudre $\begin{cases} 3x-2y = 4\\ x+y = 3\end{cases}$ par la
        méthode de votre choix.
  \item Calculer $\begin{vmatrix} 1 & -2 & 0\\ 3 & 1 & 2\\ -1 & 4 & 1
        \end{vmatrix}$.
  \item Résoudre $\begin{cases} x+y+z = 6\\ x-y = 0\\ 2x+z = 5\end{cases}$.
  \item Pour quelle valeur de $k$ le système
        $\begin{cases} 2x+ky = 1\\ 6x+9y = 4\end{cases}$
        n'a-t-il pas de solution unique ?
  \item La droite d'équation $y = ax+b$ passe par $(1\,;\,5)$ et
        $(4\,;\,14)$. Déterminer $a$ et $b$.
\end{enumerate}
\end{jemeteste}

\begin{corrige}[je me teste]
\textbf{1.} $y = 3-x$, d'où $3x-6+2x = 4$ et $x = 2$ : $(2\,;\,1)$.
\textbf{2.} Par Sarrus : $1(1)(1)+(-2)(2)(-1)+0-0-1\times4\times2
-(-2)(3)(1) = 1+4+0-0-8+6 = 3$.
\textbf{3.} $x = y$ et $z = 5-2x$ ; alors $x+x+5-2x = 6$ donne $5 = 6$ :
le système n'a aucune solution.
\textbf{4.} $D = 18-6k = 0$ pour $k = 3$.
\textbf{5.} $3a = 9$, donc $a = 3$ et $b = 2$.
\end{corrige}

% ===========================================================================
\begin{commeaubac}
\bmtsource{Exercice original au format de l'épreuve — série A \textbullet\ 5 points}
Aucun sujet de la série ne consacre d'exercice entier aux systèmes : celui-ci
est composé au format et au barème de l'épreuve.

Une école a inscrit $180$ élèves au baccalauréat, répartis entre les séries
A, S et OSE. On note $x$, $y$ et $z$ les effectifs respectifs de ces trois
séries.

\begin{enumerate}[label=\textbf{\arabic*.}, leftmargin=*, itemsep=0.28em]
  \item Traduire par une équation la phrase : « l'école a inscrit $180$
        élèves en tout ». \bareme{0,5 pt}
  \item La série A compte $20$ élèves de plus que la série OSE, et la série S
        compte autant d'élèves que les deux autres réunies. Traduire chacune
        de ces deux phrases par une équation. \bareme{1 pt}
  \item Écrire le système $(S)$ obtenu, sous forme ordonnée en $x$, $y$, $z$.
        \bareme{0,5 pt}
  \item Calculer le déterminant de $(S)$ et conclure quant au nombre de
        solutions. \bareme{1,5 pt}
  \item Résoudre $(S)$ et donner l'effectif de chaque série. \bareme{1,5 pt}
\end{enumerate}
\end{commeaubac}

\begin{corrige}[comme au baccalauréat]
\textbf{1.} $x+y+z = 180$.

\textbf{2.} « A compte $20$ de plus que OSE » : $x = z+20$, soit
$x-z = 20$. « S compte autant que les deux autres réunies » :
$y = x+z$, soit $-x+y-z = 0$.

\textbf{3.} $(S) : \begin{cases} x+y+z = 180\\ x+0\,y-z = 20\\
-x+y-z = 0\end{cases}$

\textbf{4.} Par Sarrus,
$D = (1)(0)(-1)+(1)(-1)(-1)+(1)(1)(1)-(1)(-1)(0)-(1)(1)(-1)-(1)(-1)(1)
= 0+1+1-0+1+1 = 4$. Comme $D \neq 0$, le système admet un unique triplet
solution.

\textbf{5.} La troisième équation donne $y = x+z$ ; en reportant dans la
première : $2x+2z = 180$, soit $x+z = 90$. Avec $x-z = 20$, on obtient
$x = 55$ et $z = 35$, puis $y = 90$. L'école a inscrit $55$ élèves en série
A, $90$ en série S et $35$ en série OSE. Vérification :
$55+90+35 = 180$. \checkmark
\end{corrige}

\begin{notepourlaclasse}
La méthode de Cramer à l'ordre trois est mécanique mais longue : comptez une
séance entière pour que la règle de Sarrus soit posée sans erreur de signe.
L'astuce qui fonctionne le mieux en classe est de faire recopier
systématiquement les deux colonnes supplémentaires \emph{au crayon}, quitte à
les effacer ensuite.

Le bloc d'annales mérite d'être traité avant le chapitre de statistique, et
d'être annoncé comme tel : les élèves comprennent alors que ce chapitre n'est
pas un passage obligé sans lendemain, mais l'outil d'une question qui tombe à
chaque session.

Pour une classe à grand effectif, la situation « coopérative de vanille » se
prête bien au travail de groupe : quatre élèves, quatre rôles — nommer les
inconnues, écrire les équations, résoudre, vérifier — puis rotation.
\end{notepourlaclasse}
